% Options for packages loaded elsewhere
% Options for packages loaded elsewhere
\PassOptionsToPackage{unicode}{hyperref}
\PassOptionsToPackage{hyphens}{url}
\PassOptionsToPackage{dvipsnames,svgnames,x11names}{xcolor}
%
\documentclass[
  11pt,
]{article}
\usepackage{xcolor}
\usepackage[margin=1in]{geometry}
\usepackage{amsmath,amssymb}
\setcounter{secnumdepth}{-\maxdimen} % remove section numbering
\usepackage{iftex}
\ifPDFTeX
  \usepackage[T1]{fontenc}
  \usepackage[utf8]{inputenc}
  \usepackage{textcomp} % provide euro and other symbols
\else % if luatex or xetex
  \usepackage{unicode-math} % this also loads fontspec
  \defaultfontfeatures{Scale=MatchLowercase}
  \defaultfontfeatures[\rmfamily]{Ligatures=TeX,Scale=1}
\fi
\usepackage{lmodern}
\ifPDFTeX\else
  % xetex/luatex font selection
\fi
% Use upquote if available, for straight quotes in verbatim environments
\IfFileExists{upquote.sty}{\usepackage{upquote}}{}
\IfFileExists{microtype.sty}{% use microtype if available
  \usepackage[]{microtype}
  \UseMicrotypeSet[protrusion]{basicmath} % disable protrusion for tt fonts
}{}
\makeatletter
\@ifundefined{KOMAClassName}{% if non-KOMA class
  \IfFileExists{parskip.sty}{%
    \usepackage{parskip}
  }{% else
    \setlength{\parindent}{0pt}
    \setlength{\parskip}{6pt plus 2pt minus 1pt}}
}{% if KOMA class
  \KOMAoptions{parskip=half}}
\makeatother
% Make \paragraph and \subparagraph free-standing
\makeatletter
\ifx\paragraph\undefined\else
  \let\oldparagraph\paragraph
  \renewcommand{\paragraph}{
    \@ifstar
      \xxxParagraphStar
      \xxxParagraphNoStar
  }
  \newcommand{\xxxParagraphStar}[1]{\oldparagraph*{#1}\mbox{}}
  \newcommand{\xxxParagraphNoStar}[1]{\oldparagraph{#1}\mbox{}}
\fi
\ifx\subparagraph\undefined\else
  \let\oldsubparagraph\subparagraph
  \renewcommand{\subparagraph}{
    \@ifstar
      \xxxSubParagraphStar
      \xxxSubParagraphNoStar
  }
  \newcommand{\xxxSubParagraphStar}[1]{\oldsubparagraph*{#1}\mbox{}}
  \newcommand{\xxxSubParagraphNoStar}[1]{\oldsubparagraph{#1}\mbox{}}
\fi
\makeatother


\usepackage{longtable,booktabs,array}
\newcounter{none} % for unnumbered tables
\usepackage{calc} % for calculating minipage widths
% Correct order of tables after \paragraph or \subparagraph
\usepackage{etoolbox}
\makeatletter
\patchcmd\longtable{\par}{\if@noskipsec\mbox{}\fi\par}{}{}
\makeatother
% Allow footnotes in longtable head/foot
\IfFileExists{footnotehyper.sty}{\usepackage{footnotehyper}}{\usepackage{footnote}}
\makesavenoteenv{longtable}
\usepackage{graphicx}
\makeatletter
\newsavebox\pandoc@box
\newcommand*\pandocbounded[1]{% scales image to fit in text height/width
  \sbox\pandoc@box{#1}%
  \Gscale@div\@tempa{\textheight}{\dimexpr\ht\pandoc@box+\dp\pandoc@box\relax}%
  \Gscale@div\@tempb{\linewidth}{\wd\pandoc@box}%
  \ifdim\@tempb\p@<\@tempa\p@\let\@tempa\@tempb\fi% select the smaller of both
  \ifdim\@tempa\p@<\p@\scalebox{\@tempa}{\usebox\pandoc@box}%
  \else\usebox{\pandoc@box}%
  \fi%
}
% Set default figure placement to htbp
\def\fps@figure{htbp}
\makeatother





\setlength{\emergencystretch}{3em} % prevent overfull lines

\providecommand{\tightlist}{%
  \setlength{\itemsep}{0pt}\setlength{\parskip}{0pt}}



 


\makeatletter
\@ifpackageloaded{tcolorbox}{}{\usepackage[skins,breakable]{tcolorbox}}
\@ifpackageloaded{fontawesome5}{}{\usepackage{fontawesome5}}
\definecolor{quarto-callout-color}{HTML}{909090}
\definecolor{quarto-callout-note-color}{HTML}{0758E5}
\definecolor{quarto-callout-important-color}{HTML}{CC1914}
\definecolor{quarto-callout-warning-color}{HTML}{EB9113}
\definecolor{quarto-callout-tip-color}{HTML}{00A047}
\definecolor{quarto-callout-caution-color}{HTML}{FC5300}
\definecolor{quarto-callout-color-frame}{HTML}{acacac}
\definecolor{quarto-callout-note-color-frame}{HTML}{4582ec}
\definecolor{quarto-callout-important-color-frame}{HTML}{d9534f}
\definecolor{quarto-callout-warning-color-frame}{HTML}{f0ad4e}
\definecolor{quarto-callout-tip-color-frame}{HTML}{02b875}
\definecolor{quarto-callout-caution-color-frame}{HTML}{fd7e14}
\makeatother
\makeatletter
\@ifpackageloaded{caption}{}{\usepackage{caption}}
\AtBeginDocument{%
\ifdefined\contentsname
  \renewcommand*\contentsname{Table of contents}
\else
  \newcommand\contentsname{Table of contents}
\fi
\ifdefined\listfigurename
  \renewcommand*\listfigurename{List of Figures}
\else
  \newcommand\listfigurename{List of Figures}
\fi
\ifdefined\listtablename
  \renewcommand*\listtablename{List of Tables}
\else
  \newcommand\listtablename{List of Tables}
\fi
\ifdefined\figurename
  \renewcommand*\figurename{Figure}
\else
  \newcommand\figurename{Figure}
\fi
\ifdefined\tablename
  \renewcommand*\tablename{Table}
\else
  \newcommand\tablename{Table}
\fi
}
\@ifpackageloaded{float}{}{\usepackage{float}}
\floatstyle{ruled}
\@ifundefined{c@chapter}{\newfloat{codelisting}{h}{lop}}{\newfloat{codelisting}{h}{lop}[chapter]}
\floatname{codelisting}{Listing}
\newcommand*\listoflistings{\listof{codelisting}{List of Listings}}
\usepackage{amsthm}
\theoremstyle{plain}
\newtheorem{theorem}{Theorem}[section]
\theoremstyle{plain}
\newtheorem{lemma}{Lemma}[section]
\theoremstyle{remark}
\AtBeginDocument{\renewcommand*{\proofname}{Proof}}
\newtheorem*{remark}{Remark}
\newtheorem*{solution}{Solution}
\newtheorem{refremark}{Remark}[section]
\newtheorem{refsolution}{Solution}[section]
\makeatother
\makeatletter
\makeatother
\makeatletter
\@ifpackageloaded{caption}{}{\usepackage{caption}}
\@ifpackageloaded{subcaption}{}{\usepackage{subcaption}}
\makeatother
\usepackage{bookmark}
\IfFileExists{xurl.sty}{\usepackage{xurl}}{} % add URL line breaks if available
\urlstyle{same}
% fallback for those not using the hyperref driver hyperxmp:
\makeatletter
\@ifundefined{xmpquote}{\newcommand{\xmpquote}[1]{#1}}{}
\makeatother
\hypersetup{
  pdftitle={Dvir's proof of the finite field Kakeya conjecture},
  pdfauthor={Hung-Hsun Hans Yu},
  colorlinks=true,
  linkcolor={blue},
  filecolor={Maroon},
  citecolor={Blue},
  urlcolor={Blue},
  pdfcreator={LaTeX via pandoc}}


\title{Dvir's proof of the finite field Kakeya conjecture}
\author{Hung-Hsun Hans Yu}
\date{2026-10-16}
\begin{document}
\maketitle

\renewcommand*\contentsname{Table of contents}
{
\setcounter{tocdepth}{3}
\tableofcontents
}

\providecommand{\F}{\mathbb{F}}
\providecommand{\R}{\mathbb{R}}
\providecommand{\C}{\mathbb{C}}
\providecommand{\Z}{\mathbb{Z}}
\providecommand{\N}{\mathbb{N}}
\providecommand{\Q}{\mathbb{Q}}
\providecommand{\E}{\mathbb{E}}
\providecommand{\P}{\mathbb{P}}
\providecommand{\eps}{\varepsilon}
\providecommand{\abs}[1]{\left|#1\right|}
\providecommand{\norm}[1]{\left\lVert#1\right\rVert}
\providecommand{\set}[1]{\left\{#1\right\}}
\providecommand{\floor}[1]{\left\lfloor#1\right\rfloor}
\providecommand{\ceil}[1]{\left\lceil#1\right\rceil}
\providecommand{\poly}{\operatorname{poly}}

A \emph{Kakeya set} in \(\mathbb{F}_q^n\) is a set \(K\) that contains a
line in every direction: for every
\(v \in \mathbb{F}_q^n \setminus \left\{0\right\}\) there is a point
\(x\) with \(\left\{x + tv : t \in \mathbb{F}_q\right\} \subseteq K\).
Wolff asked whether every such set has size at least
\(c_n q^n\),\footnote{Wolff raised the question in 1999 as a finite
  field model of the Euclidean Kakeya conjecture.} and Dvir answered the
question with a remarkably short argument.

\subsection{Two facts about
polynomials}\label{two-facts-about-polynomials}

\begin{lemma}[Interpolation]\protect\hypertarget{lem-interpolation}{}\label{lem-interpolation}

If \(E \subseteq \mathbb{F}^n\) is finite and
\(\left|E\right| < \binom{d+n}{n}\), then some nonzero polynomial
\(P \in \mathbb{F}[x_1, \dots, x_n]\) of degree at most \(d\) vanishes
on \(E\).

\end{lemma}

\begin{proof}
The polynomials of degree at most \(d\) form a vector space of dimension

\begin{equation}\protect\phantomsection\label{eq-dimension}{
\dim \mathbb{F}[x_1, \dots, x_n]_{\le d} = \binom{d+n}{n},
}\end{equation}

since the right side counts monomials of degree at most \(d\).
Evaluation at the points of \(E\) is a linear map from this space to
\(\mathbb{F}^E\). By Equation~\ref{eq-dimension} the domain has larger
dimension than the codomain, so the kernel is nontrivial.
\end{proof}

\begin{lemma}[]\protect\hypertarget{lem-nonvanishing}{}\label{lem-nonvanishing}

If \(Q \in \mathbb{F}_q[x_1, \dots, x_n]\) is nonzero and has degree
less than \(q\) in each variable, then \(Q\) does not vanish on all of
\(\mathbb{F}_q^n\).

\end{lemma}

\begin{proof}
Induct on \(n\). For \(n = 1\), a nonzero polynomial of degree less than
\(q\) has fewer than \(q\) roots. For \(n > 1\), write
\(Q = \sum_{i < q} Q_i(x_1, \dots, x_{n-1})\, x_n^i\) with some
\(Q_i \ne 0\). By induction there is \(a \in \mathbb{F}_q^{n-1}\) with
\(Q_i(a) \ne 0\), so \(Q(a, x_n)\) is a nonzero polynomial in \(x_n\) of
degree less than \(q\), and it has a non-root.
\end{proof}

\subsection{The theorem}\label{the-theorem}

\begin{theorem}[Dvir,
2009]\protect\hypertarget{thm-dvir}{}\label{thm-dvir}

Every Kakeya set \(K \subseteq \mathbb{F}_q^n\) satisfies \[
\left|K\right| \ge \binom{q+n-1}{n} \ge \frac{q^n}{n!}.
\]

\end{theorem}

\begin{proof}
Suppose \(\left|K\right| < \binom{q+n-1}{n}\). Applying
Lemma~\ref{lem-interpolation} with \(d = q - 1\) gives a nonzero
polynomial \(P\) of degree at most \(q - 1\) that vanishes on \(K\). Let
\(d = \deg P\) and let \(P_d\) be the homogeneous part of \(P\) of
degree \(d\). Since \(K\) is nonempty, \(P\) is not a nonzero constant,
so \(d \ge 1\).

Fix a direction \(v \ne 0\) and choose \(x\) with \(x + tv \in K\) for
all \(t \in \mathbb{F}_q\). The polynomial \(g(t) = P(x + tv)\) has
degree at most \(d \le q - 1\) and vanishes at all \(q\) elements of
\(\mathbb{F}_q\), so \(g \equiv 0\). Its coefficient of \(t^d\) is
\(P_d(v)\), hence \(P_d(v) = 0\). Since \(P_d\) is homogeneous of degree
\(d \ge 1\), also \(P_d(0) = 0\).

So \(P_d\) vanishes on all of \(\mathbb{F}_q^n\), yet it is nonzero with
degree less than \(q\) in each variable. This contradicts
Lemma~\ref{lem-nonvanishing}.
\end{proof}

\begin{tcolorbox}[enhanced jigsaw, arc=.35mm, bottomrule=.15mm, bottomtitle=1mm, breakable, colback=white, colbacktitle=quarto-callout-note-color!10!white, colframe=quarto-callout-note-color-frame, coltitle=black, left=2mm, leftrule=.75mm, opacityback=0, opacitybacktitle=0.6, rightrule=.15mm, title=\textcolor{quarto-callout-note-color}{\faInfo}\hspace{0.5em}{Remark}, titlerule=0mm, toprule=.15mm, toptitle=1mm]

Using polynomials that vanish to high multiplicity, Dvir, Kopparty,
Saraf, and Sudan improved the constant in Theorem~\ref{thm-dvir} from
\(1/n!\) to \(2^{-n}\).

\end{tcolorbox}




\end{document}
